\documentclass{amsart}
% For dealing with references we use the comment environment
\usepackage{verbatim}
\newenvironment{reference}{\comment}{\endcomment}
%\newenvironment{reference}{}{}
\newenvironment{slogan}{\comment}{\endcomment}
\newenvironment{history}{\comment}{\endcomment}

% For commutative diagrams we use Xy-pic
\usepackage[all]{xy}

% We use 2cell for 2-commutative diagrams.
\xyoption{2cell}
\UseAllTwocells

% We use multicol for the list of chapters between chapters
\usepackage{multicol}

% This is generally recommended for better output
\usepackage{lmodern}
\usepackage[T1]{fontenc}

% For cross-file-references

% Package for hypertext links:
\usepackage{hyperref}

% For any local file, say "hello.tex" you want to link to please

% Theorem environments.
%
\theoremstyle{plain}
\newtheorem{theorem}[subsection]{Theorem}
\newtheorem{proposition}[subsection]{Proposition}
\newtheorem{lemma}[subsection]{Lemma}

\theoremstyle{definition}
\newtheorem{definition}[subsection]{Definition}
\newtheorem{example}[subsection]{Example}
\newtheorem{exercise}[subsection]{Exercise}
\newtheorem{situation}[subsection]{Situation}

\theoremstyle{remark}
\newtheorem{remark}[subsection]{Remark}
\newtheorem{remarks}[subsection]{Remarks}

\numberwithin{equation}{subsection}

% Macros
%
\def\lim{\mathop{\mathrm{lim}}\nolimits}
\def\colim{\mathop{\mathrm{colim}}\nolimits}
\def\Spec{\mathop{\mathrm{Spec}}}
\def\Hom{\mathop{\mathrm{Hom}}\nolimits}
\def\Ext{\mathop{\mathrm{Ext}}\nolimits}
\def\SheafHom{\mathop{\mathcal{H}\!\mathit{om}}\nolimits}
\def\SheafExt{\mathop{\mathcal{E}\!\mathit{xt}}\nolimits}
\def\Sch{\mathit{Sch}}
\def\Mor{\mathop{\mathrm{Mor}}\nolimits}
\def\Ob{\mathop{\mathrm{Ob}}\nolimits}
\def\Sh{\mathop{\mathit{Sh}}\nolimits}
\def\NL{\mathop{N\!L}\nolimits}
\def\CH{\mathop{\mathrm{CH}}\nolimits}
\def\proetale{{pro\text{-}\acute{e}tale}}
\def\etale{{\acute{e}tale}}
\def\QCoh{\mathit{QCoh}}
\def\Ker{\mathop{\mathrm{Ker}}}
\def\Im{\mathop{\mathrm{Im}}}
\def\Coker{\mathop{\mathrm{Coker}}}
\def\Coim{\mathop{\mathrm{Coim}}}

% Boxtimes
%
\DeclareMathSymbol{\boxtimes}{\mathbin}{AMSa}{"02}

%
% Macros for moduli stacks/spaces
%
\def\QCohstack{\mathcal{QC}\!\mathit{oh}}
\def\Cohstack{\mathcal{C}\!\mathit{oh}}
\def\Spacesstack{\mathcal{S}\!\mathit{paces}}
\def\Quotfunctor{\mathrm{Quot}}
\def\Hilbfunctor{\mathrm{Hilb}}
\def\Curvesstack{\mathcal{C}\!\mathit{urves}}
\def\Polarizedstack{\mathcal{P}\!\mathit{olarized}}
\def\Complexesstack{\mathcal{C}\!\mathit{omplexes}}
% \Pic is the operator that assigns to X its picard group, usage \Pic(X)
% \Picardstack_{X/B} denotes the Picard stack of X over B
% \Picardfunctor_{X/B} denotes the Picard functor of X over B
\def\Pic{\mathop{\mathrm{Pic}}\nolimits}
\def\Picardstack{\mathcal{P}\!\mathit{ic}}
\def\Picardfunctor{\mathrm{Pic}}
\def\Deformationcategory{\mathcal{D}\!\mathit{ef}}

\setlength{\emergencystretch}{3em}
\begin{document}
\title[Stacks Project, Tag 01A4]{707 --- Chain homotopies lift to prescribed simplicial homotopies in abelian categories}
\maketitle
\noindent Copyright (C) 2005--2025 Johan de Jong. Stacks Project authors. Source: \href{https://stacks.math.columbia.edu/tag/01A4}{Tag 01A4}.

\noindent Editorial difficulty note (not author text). Difficulty H2: The author compares a chain cylinder with the simplicial interval by constructing explicit maps in both directions and computing their sign-sensitive interaction. Recovering the prescribed normalized homotopy is a substantial converse to ordinary homotopy normalization..

\noindent Editorial provenance note (not author text). History: excerpt from revision \texttt{a0444\allowbreak 6e57e\allowbreak c1fbc\allowbreak 252a8\allowbreak 71afc\allowbreak ec775\allowbreak 2fb28\allowbreak 07b14}, simplicial.tex, lines 5289--5500. Reference presentation was adapted; original-author context and core support blocks were retained or added. The mathematical wording is unchanged. Licensed under GNU FDL 1.2 or later; no invariant sections or cover texts. See ../licenses/COPYING. Context blocks reproduce author definitions and supporting results; they are not additional counted problems.

\begin{definition}
\label{definition-homotopy}
Let $\mathcal{C}$ be a category having finite coproducts.
Suppose that $U$ and $V$ are two simplicial objects of $\mathcal{C}$.
Let $a, b : U \to V$ be two morphisms.
\begin{enumerate}
\item We say a morphism
$$
h : U \times \Delta[1] \longrightarrow V
$$
is a {\it homotopy from $a$ to $b$} if $a = h \circ e_0$ and
$b = h \circ e_1$.
\item We say the morphisms $a$ and $b$ are {\it homotopic} or are
{\it in the same homotopy class}
if there exists a sequence of morphisms $a = a_0, a_1, \ldots, a_n = b$
from $U$ to $V$ such that for each $i = 1, \ldots, n$ there either exists
a homotopy from $a_{i - 1}$ to $a_i$ or there exists a homotopy
from $a_i$ to $a_{i - 1}$.
\end{enumerate}
\end{definition}

\noindent
There is a second chain complex we can associate to
a simplicial object of $\mathcal{A}$. Recall that by
Lemma \href{https://stacks.math.columbia.edu/tag/017V}{lemma splitting abelian category (Tag 017V)}
any simplicial object $U$ of $\mathcal{A}$
is canonically split with
$N(U_m) = \bigcap_{i = 0}^{m - 1} \Ker(d^m_i)$.
We define the {\it normalized chain complex $N(U)$}
to be the chain complex
$$
\ldots \to N(U_2) \to N(U_1) \to N(U_0) \to 0 \to 0 \to \ldots
$$
with boundary map $d_n : N(U_n) \to N(U_{n - 1})$ given
by the restriction of $(-1)^nd^n_n$ to the direct summand
$N(U_n)$ of $U_n$. Note that Lemma \href{https://stacks.math.columbia.edu/tag/017X}{lemma N d in N (Tag 017X)}
implies that $d^n_n(N(U_n)) \subset N(U_{n - 1})$.
It is a complex because
$d^n_n \circ d^{n + 1}_{n + 1} = d^n_n \circ d^{n + 1}_n$
and $d^{n + 1}_n$ is zero on $N(U_{n + 1})$ by definition.
Thus we obtain a second functor
$$
N : \text{Simp}(\mathcal{A}) \longrightarrow \text{Ch}_{\geq 0}(\mathcal{A}).
$$


\begin{lemma}
\label{lemma-represent-homotopy}
Let $\mathcal{A}$ be an abelian category. Let $A$ be a chain complex.
Consider the covariant functor
$$
B \longmapsto
\{
(a, b, h)
\mid
a, b : A \to B\text{ and }h\text{ a homotopy between }a, b
\}
$$
There exists a chain complex $\diamond A$
such that $\Mor_{\text{Ch}(\mathcal{A})}(\diamond A, -)$
is isomorphic to the displayed functor.
The construction $A \mapsto \diamond A$ is functorial.
\end{lemma}

\begin{proof}
We set $\diamond A_n = A_n \oplus A_n \oplus A_{n - 1}$,
and we define $d_{\diamond A, n}$ by the matrix
$$
d_{\diamond A, n}
=
\left(
\begin{matrix}
d_{A, n} & 0 & \text{id}_{A_{n - 1}} \\
0 & d_{A, n} & -\text{id}_{A_{n - 1}} \\
0 & 0 & -d_{A, n - 1}
\end{matrix}
\right) :
A_n \oplus A_n \oplus A_{n - 1} \to
A_{n - 1} \oplus A_{n - 1} \oplus A_{n - 2}
$$
If $\mathcal{A}$ is the category of abelian groups,
and $(x, y, z) \in A_n \oplus A_n \oplus A_{n - 1}$
then $d_{\diamond A, n}(x, y, z) = (d_n(x) + z, d_n(y) - z, -d_{n - 1}(z))$.
It is easy to verify that $d^2 = 0$. Clearly,
there are two maps $\diamond a, \diamond b : A \to \diamond A$
(first summand and second summand),
and a map $\diamond A \to A[-1]$ which give
a short exact sequence
$$
0 \to
A \oplus A \to
\diamond A \to
A[-1] \to
0
$$
which is termwise split. Moreover, there is a sequence
of maps $\diamond h_n : A_n \to \diamond A_{n + 1}$, namely
the identity from $A_n$ to the summand $A_n$ of $\diamond A_{n + 1}$,
such that $\diamond h$ is a homotopy between $\diamond a$ and $\diamond b$.

\medskip\noindent
We conclude that any morphism $f : \diamond A \to B$
gives rise to a triple $(a, b, h)$ by setting $a = f \circ \diamond a$,
$b = f \circ \diamond b$ and $h_n = f_{n + 1} \circ \diamond h_n$.
Conversely, given a triple $(a, b, h)$ we get a morphism
$f : \diamond A \to B$ by taking
$$
f_n = (a_n, b_n, h_{n - 1}).
$$
To see that this is a morphism of chain complexes you have
to do a calculation. We only do this in case $\mathcal{A}$
is the category of abelian groups:
Say $(x, y, z) \in \diamond A_n = A_n \oplus A_n \oplus A_{n - 1}$.
Then
\begin{eqnarray*}
f_{n - 1}(d_n(x, y, z)) & = &
f_{n - 1}(d_n(x) + z, d_n(y) - z, -d_{n - 1}(z)) \\
& = &
a_n(d_n(x)) + a_n(z) + b_n(d_n(y)) - b_n(z) - h_{n - 2}(d_{n - 1}(z))
\end{eqnarray*}
and
\begin{eqnarray*}
d_n(f_n(x, y, z) & = &
d_n(a_n(x) + b_n(y) + h_{n - 1}(z)) \\
& = &
d_n(a_n(x)) + d_n(b_n(y)) + d_n(h_{n - 1}(z))
\end{eqnarray*}
which are the same by definition of a homotopy.
\end{proof}


\begin{lemma}
\label{lemma-map-into-diamond}
Let $\mathcal{A}$ be an abelian category.
Let
$$
0 \to A \oplus A \to B \to C \to 0
$$
be a short exact sequence of chain complexes of $\mathcal{A}$.
Suppose given in addition morphisms $s_n : C_n \to B_n$
splitting the associated short exact sequence in degree $n$.
Let $\delta(s) : C \to (A \oplus A)[-1] = A[-1] \oplus A[-1]$
be the associated morphism of complexes, see
Homology, Lemma \href{https://stacks.math.columbia.edu/tag/011D}{lemma ses termwise split (Tag 011D)}.
If $\delta(s)$ factors through the morphism
$(1, -1) : A[-1] \to A[-1] \oplus A[-1]$, then
there is a unique morphism $B \to \diamond A$
fitting into a commutative diagram
$$
\xymatrix{
0 \ar[r] &
A \oplus A \ar[d] \ar[r] &
B \ar[r] \ar[d] &
C \ar[d] \ar[r] &
0 \\
0 \ar[r] &
A \oplus A \ar[r] &
\diamond A \ar[r] &
A[-1] \ar[r] &
0
}
$$
where the vertical maps are compatible with the splittings
$s_n$ and the splittings of $\diamond A_n \to A[-1]_n$
as well.
\end{lemma}

\begin{proof}
Denote $(p_n, q_n) : B_n \to A_n \oplus A_n$ the morphism $\pi_n$
of Homology, Lemma \href{https://stacks.math.columbia.edu/tag/011D}{lemma ses termwise split (Tag 011D)}.
Also write $(a, b) : A \oplus A \to B$, and
$r : B \to C$ for the maps in the short exact sequence.
Write the factorization of $\delta(s)$ as
$\delta(s) = (1, -1) \circ f$. This means that
$p_{n - 1} \circ d_{B, n} \circ s_n = f_n$, and
$q_{n - 1} \circ d_{B, n} \circ s_n = - f_n$, and
Set $B_n \to \diamond A_n = A_n \oplus A_n \oplus A_{n - 1}$
equal to $(p_n, q_n, f_n \circ r_n)$.

\medskip\noindent
Now we have to check that this actually defines a morphism
of complexes. We will only do this in the case of abelian groups.
Pick $x \in B_n$. Then $x = a_n(x_1) + b_n(x_2) + s_n(x_3)$
and it suffices to show that our definition commutes
with differential for each term separately. For the term
$a_n(x_1)$ we have $(p_n, q_n, f_n \circ r_n)(a_n(x_1)) =
(x_1, 0, 0)$ and the result is obvious. Similarly for
the term $b_n(x_2)$. For the term $s_n(x_3)$ we have
\begin{eqnarray*}
(p_n, q_n, f_n \circ r_n)(d_n(s_n(x_3))) & = &
(p_n, q_n, f_n \circ r_n)( \\
& & \ \ \ \ \ a_n(f_n(x_3)) - b_n(f_n(x_3)) + s_n(d_n(x_3))) \\
& = &
(f_n(x_3), -f_n(x_3), f_n(d_n(x_3)))
\end{eqnarray*}
by definition of $f_n$. And
\begin{eqnarray*}
d_n(p_n, q_n, f_n \circ r_n)(s_n(x_3)) & = & d_n(0, 0, f_n(x_3)) \\
& = &
(f_n(x_3), - f_n(x_3), d_{A[-1], n}(f_n(x_3)))
\end{eqnarray*}
The result follows as $f$ is a morphism of complexes.
\end{proof}


\begin{lemma}
\label{lemma-backwards-homotopy}
Let $\mathcal{A}$ be an abelian category.
Let $U$, $V$ be simplicial objects of $\mathcal{A}$.
Let $a, b : U \to V$ be a pair of morphisms.
Assume the corresponding maps of chain complexes
$N(a), N(b) : N(U) \to N(V)$ are homotopic by
a homotopy $\{N_n : N(U)_n \to N(V)_{n + 1}\}$.
Then there exists a homotopy from $a$ to $b$ as in
Definition \ref{definition-homotopy}. Moreover, one can choose the
homotopy $h : U \times \Delta[1] \to V$ such that
$N_n = N(h)_n$ where $N(h)$ is the homotopy coming
from $h$ as in Section \href{https://stacks.math.columbia.edu/tag/019Q}{section homotopy abelian (Tag 019Q)}.
\end{lemma}

\begin{proof}
Let $(\diamond N(U), \diamond a, \diamond b, \diamond h)$
be as in Lemma \ref{lemma-represent-homotopy} and its proof. By that lemma
there exists a morphism $\diamond N(U) \to N(V)$ representing
the triple $(N(a), N(b), \{N_n\})$. We will show
there exists a morphism
$\psi : N(U \times \Delta[1]) \to \diamond{N(U)}$
such that $\diamond a = \psi \circ N(e_0)$, and
$\diamond b = \psi \circ N(e_1)$. Moreover, we will
show that the homotopy between $N(e_0), N(e_1) :
N(U) \to N(U \times \Delta[1])$ coming from
(\href{https://stacks.math.columbia.edu/tag/019R}{equation homotopy to homotopy (Tag 019R)}) and
Lemma \href{https://stacks.math.columbia.edu/tag/019S}{lemma homotopy s N (Tag 019S)} with
$h = \text{id}_{U \times \Delta[1]}$ is mapped via
$\psi$ to the canonical homotopy $\diamond h$ between the two
maps $\diamond a, \diamond b : N(U) \to \diamond{N(U)}$. Certainly this
will imply the lemma.

\medskip\noindent
Note that $N : \text{Simp}(\mathcal{A}) \to \text{Ch}_{\geq 0}(\mathcal{A})$
as a functor is a direct summand of the functor
$s : \text{Simp}(\mathcal{A}) \to \text{Ch}_{\geq 0}(\mathcal{A})$.
Also, the functor $\diamond$ is compatible with direct sums.
Thus it suffices instead to construct a morphism
$\Psi : s(U \times \Delta[1]) \to \diamond{s(U)}$ with the
corresponding properties. This is what we do below.

\medskip\noindent
By Definition \ref{definition-homotopy}
the morphisms $e_0 : U \to U \times \Delta[1]$
and $e_1 : U \to U \times \Delta[1]$ are homotopic
with homotopy $\text{id}_{U \times \Delta[1]}$.
By Lemma \href{https://stacks.math.columbia.edu/tag/019S}{lemma homotopy s N (Tag 019S)}
we get an explicit homotopy
$\{h_n : s(U)_n \to s(U \times \Delta[1])_{n + 1}\}$
between the morphisms
of chain complexes $s(e_0) : s(U) \to s(U \times \Delta[1])$
and $s(e_1) : s(U) \to s(U \times \Delta[1])$. By
Lemma \ref{lemma-map-into-diamond} above we get a corresponding morphism
$$
\Phi : \diamond{s(U)} \to s(U \times \Delta[1])
$$
According to the construction, $\Phi_n$ restricted to the summand
$s(U)[-1]_n = s(U)_{n - 1}$ of $\diamond{s(U)}_n$
is equal to $h_{n - 1}$. And
$$
h_{n - 1} = \sum\nolimits_{i = 0}^{n - 1}
(-1)^{i + 1} s^n_i \cdot \alpha^n_{i + 1} :
U_{n - 1} \to \bigoplus\nolimits_j U_n \cdot \alpha^n_j.
$$
with obvious notation.

\medskip\noindent
On the other hand, the morphisms $e_i : U \to U \times \Delta[1]$ induce
a morphism $(e_0, e_1) : U \oplus U \to U \times \Delta[1]$.
Denote $W$ the cokernel. Note that, if we write
$(U \times \Delta[1])_n = \bigoplus_{\alpha : [n] \to [1]} U_n \cdot \alpha$,
then we may identify
$W_n = \bigoplus_{i = 1}^n U_n \cdot \alpha^n_i$ with
$\alpha^n_i$ as in Section \href{https://stacks.math.columbia.edu/tag/019J}{section homotopy (Tag 019J)}.
We have a commutative diagram
$$
\xymatrix{
0 \ar[r] &
U \oplus U \ar[rd]_{(1, 1)} \ar[r] &
U \times \Delta[1] \ar[d]^\pi \ar[r] &
W \ar[r] &
0 \\
& & U & &
}
$$
This implies we have a similar commutative diagram after
applying the functor $s$. Next, we choose the splittings
$\sigma_n : s(W)_n \to s(U \times \Delta[1])_n$ by mapping
the summand $U_n \cdot \alpha^n_i \subset W_n$ via $(-1, 1)$
to the summands $U_n \cdot \alpha^n_0 \oplus U_n \cdot \alpha^n_i
\subset (U \times \Delta[1])_n$. Note that $s(\pi)_n \circ \sigma_n = 0$.
It follows that $(1, 1) \circ \delta(\sigma)_n = 0$.
Hence $\delta(\sigma)$ factors as in
Lemma \ref{lemma-map-into-diamond}. By that lemma
we obtain a canonical morphism
$\Psi : s(U \times \Delta[1]) \to \diamond{s(U)}$.

\medskip\noindent
To compute $\Psi$ we first compute the morphism
$\delta(\sigma) : s(W) \to s(U)[-1] \oplus s(U)[-1]$.
According to Homology, Lemma \href{https://stacks.math.columbia.edu/tag/011D}{lemma ses termwise split (Tag 011D)}
and its proof,
to do this we have compute
$$
d_{s(U \times \delta[1]), n} \circ \sigma_n
-
\sigma_{n - 1} \circ d_{s(W), n}
$$
and write it as a morphism into
$U_{n - 1} \cdot \alpha^{n - 1}_0 \oplus U_{n - 1} \cdot \alpha^{n - 1}_n$.
We only do this in case $\mathcal{A}$ is the category of abelian
groups. We use the short hand notation $x_{\alpha}$ for $x \in U_n$
to denote the element $x$ in the summand $U_n \cdot \alpha$
of $(U \times \Delta[1])_n$. Recall that
$$
d_{s(U \times \delta[1]), n}
=
\sum\nolimits_{i = 0}^n (-1)^i d^n_i
$$
where $d^n_i$ maps the summand $U_n \cdot \alpha$
to the summand $U_{n - 1} \cdot (\alpha \circ \delta^n_i)$
via the morphism $d^n_i$ of the simplicial object $U$.
In terms of the notation above this means
$$
d_{s(U \times \delta[1]), n}(x_\alpha) =
\sum\nolimits_{i = 0}^n (-1)^i (d^n_i(x))_{\alpha \circ \delta^n_i}
$$
Starting with $x_\alpha \in W_n$, in other words $\alpha = \alpha^n_j$
for some $j \in \{1, \ldots, n\}$, we see that
$\sigma_n(x_\alpha) = x_\alpha - x_{\alpha^n_0}$ and
hence
$$
(d_{s(U \times \delta[1]), n} \circ \sigma_n)(x_\alpha)
=
\sum\nolimits_{i = 0}^n (-1)^i (d^n_i(x))_{\alpha \circ \delta^n_i}
-
\sum\nolimits_{i = 0}^n (-1)^i (d^n_i(x))_{\alpha^n_0 \circ \delta^n_i}
$$
To compute $d_{s(W), n}(x_\alpha)$, we have to omit
all terms where $\alpha \circ \delta^n_i = \alpha^{n - 1}_0, \alpha^{n - 1}_n$.
Hence we get
$$
\begin{matrix}
(\sigma_{n - 1} \circ d_{s(W), n})(x_\alpha) = \\
\sum\nolimits_{i = 0, \ldots, n\text{ and }
\alpha \circ \delta^n_i \not = \alpha^{n - 1}_0\text{ or }\alpha^{n - 1}_n}
\Big((-1)^i (d^n_i(x))_{\alpha \circ \delta^n_i}
-
(-1)^i (d^n_i(x))_{\alpha^{n - 1}_0}
\Big)
\end{matrix}
$$
Clearly the difference of the two terms is the sum
$$
\sum\nolimits_{i = 0, \ldots, n\text{ and }
\alpha \circ \delta^n_i = \alpha^{n - 1}_0\text{ or }\alpha^{n - 1}_n}
\Big((-1)^i (d^n_i(x))_{\alpha \circ \delta^n_i}
-
(-1)^i (d^n_i(x))_{\alpha^{n - 1}_0}
\Big)
$$
Of course, if $\alpha \circ \delta^n_i = \alpha^{n - 1}_0$
then the term drops out. Recall that $\alpha = \alpha^n_j$
for some $j \in \{1, \ldots, n\}$. The only way
$\alpha^n_j \circ \delta^n_i = \alpha^{n - 1}_n$
is if $j = n$ and $i = n$. Thus we actually get
$0$ unless $j = n$ and in that case we get
$(-1)^n(d^n_n(x))_{\alpha^{n - 1}_n} - (-1)^n(d^n_n(x))_{\alpha^{n - 1}_0}$.
In other words, we conclude the morphism
$$
\delta(\sigma)_n :
W_n
\to
(s(U)[-1] \oplus s(U)[-1])_n = U_{n - 1} \oplus U_{n - 1}
$$
is zero on all summands except $U_n \cdot \alpha^n_n$
and on that summand it is equal to $((-1)^nd^n_n, -(-1)^nd^n_n)$.
(Namely, the first summand of the two corresponds to the factor
with $\alpha^{n - 1}_n$ because that is the map $[n - 1] \to [1]$
which maps everybody to $0$, and hence corresponds to $e_0$.)

\medskip\noindent
We obtain a canonical diagram
$$
\xymatrix{
0 \ar[r] &
s(U) \oplus s(U) \ar[r] \ar[d] &
\diamond{s(U)} \ar[r] \ar[d]^{\Phi}&
s(U)[-1] \ar[r] \ar[d] &
0 \\
0 \ar[r] &
s(U) \oplus s(U) \ar[r] \ar[d] &
s(U \times \Delta[1]) \ar[r] \ar[d]^\Psi &
s(W) \ar[r] \ar[d] &
0 \\
0 \ar[r] &
s(U) \oplus s(U) \ar[r] &
\diamond{s(U)} \ar[r] &
s(U)[-1] \ar[r] &
0
}
$$
We claim that $\Phi \circ \Psi$ is the identity.
To see this it is enough to prove that the composition
of $\Phi$ and $\delta(\sigma)$ as a map
$s(U)[-1] \to s(W) \to s(U)[-1] \oplus s(U)[-1]$ is the
identity in the first factor and minus identity in the second.
By the computations above it is
$((-1)^nd^n_0, -(-1)^nd^n_0) \circ (-1)^n s^n_n = (1, -1)$
as desired.
\end{proof}
\end{document}
